%% toc: P2 12 12.5,12.6
\section{Proofs for Cut, the subformula property, and the translations between natural deduction and sequents}
\label{sec:annex-p02-w02b-p02-009-cut-translations}

Notation is that of \cref{def:p02-w02b-p02-009-setting,def:p02-w02b-p02-009-validity}.
Two facts about multisets are used throughout:
\[
 \operatorname{supp}(\Pi,\Sigma)
 =\operatorname{supp}(\Pi)\cup\operatorname{supp}(\Sigma),
 \qquad
 \operatorname{supp}(\Pi,\varphi,\varphi)=\operatorname{supp}(\Pi,\varphi).
\]
The truth clauses are those of the valuation definition in
\cref{sec:prop-syntax-semantics-02}: \(v\not\entails\bottom\);
\(v\entails\neg\varphi\) iff \(v\not\entails\varphi\);
\(v\entails\varphi\land\psi\) iff \(v\entails\varphi\) and \(v\entails\psi\);
\(v\entails\varphi\lor\psi\) iff \(v\entails\varphi\) or \(v\entails\psi\);
\(v\not\entails\varphi\to\psi\) iff \(v\entails\varphi\) and
\(v\not\entails\psi\).

\begin{annexitem}{annex:p02-w02b-p02-009-lk-rule-soundness}{thm:p02-w02b-p02-009-lk-rule-soundness}{Every rule of the sequent calculus reflects falsification}
\begin{proof}
\emph{Initial sequents.}  In an identity sequent
\(\Gamma,\varphi\Rightarrow\varphi,\Delta\) the formula \(\varphi\) lies in
the support of both sides.  A valuation \(v\) with \(v\entails\varphi\) makes
a succedent formula true, and one with \(v\not\entails\varphi\) makes an
antecedent formula false; in neither case does \(v\) falsify the sequent.  In
a \(\bottom\)-sequent \(\Gamma,\bottom\Rightarrow\Delta\) the antecedent
support contains \(\bottom\), which no valuation makes true.

In each rule case below, \(v\) is a valuation falsifying the conclusion, so
every member of the antecedent support of the conclusion is true under \(v\)
and every member of its succedent support is false under \(v\).

\(W_L\) (premise \(\Gamma\Rightarrow\Delta\), conclusion
\(\Gamma,\varphi\Rightarrow\Delta\)): since
\(\operatorname{supp}(\Gamma)\subseteq\operatorname{supp}(\Gamma,\varphi)\),
every member of \(\operatorname{supp}(\Gamma)\) is true and every member of
\(\operatorname{supp}(\Delta)\) is false, so \(v\) falsifies the premise.
\(W_R\) is the same with
\(\operatorname{supp}(\Delta)\subseteq\operatorname{supp}(\Delta,\varphi)\).

\(C_L\) and \(C_R\): premise and conclusion have the same supports on each
side, so \(v\) falsifies the premise.

\(\neg L\) (premise \(\Gamma\Rightarrow\Delta,\varphi\), conclusion
\(\Gamma,\neg\varphi\Rightarrow\Delta\)): \(v\entails\neg\varphi\), so
\(v\not\entails\varphi\).  The members of \(\operatorname{supp}(\Gamma)\) are
true and those of \(\operatorname{supp}(\Delta)\cup\{\varphi\}\) false, so
\(v\) falsifies the premise.

\(\neg R\) (premise \(\Gamma,\varphi\Rightarrow\Delta\), conclusion
\(\Gamma\Rightarrow\Delta,\neg\varphi\)): \(v\not\entails\neg\varphi\), so
\(v\entails\varphi\); the members of
\(\operatorname{supp}(\Gamma)\cup\{\varphi\}\) are true and those of
\(\operatorname{supp}(\Delta)\) false.

\(\land L\) (premise \(\Gamma,\varphi,\psi\Rightarrow\Delta\), conclusion
\(\Gamma,\varphi\land\psi\Rightarrow\Delta\)): \(v\entails\varphi\land\psi\),
so \(v\entails\varphi\) and \(v\entails\psi\); \(v\) falsifies the premise.

\(\land R\) (premises \(\Gamma\Rightarrow\Delta,\varphi\) and
\(\Gamma\Rightarrow\Delta,\psi\), conclusion
\(\Gamma\Rightarrow\Delta,\varphi\land\psi\)):
\(v\not\entails\varphi\land\psi\), so \(v\not\entails\varphi\) or
\(v\not\entails\psi\).  In the first case \(v\) falsifies the left premise,
in the second the right premise.

\(\lor L\) (premises \(\Gamma,\varphi\Rightarrow\Delta\) and
\(\Gamma,\psi\Rightarrow\Delta\), conclusion
\(\Gamma,\varphi\lor\psi\Rightarrow\Delta\)): \(v\entails\varphi\lor\psi\),
so \(v\entails\varphi\) or \(v\entails\psi\); accordingly \(v\) falsifies the
left or the right premise.

\(\lor R\) (premise \(\Gamma\Rightarrow\Delta,\varphi,\psi\), conclusion
\(\Gamma\Rightarrow\Delta,\varphi\lor\psi\)):
\(v\not\entails\varphi\lor\psi\), so \(v\not\entails\varphi\) and
\(v\not\entails\psi\); \(v\) falsifies the premise.

\(\to L\) (premises \(\Gamma\Rightarrow\Delta,\varphi\) and
\(\Gamma,\psi\Rightarrow\Delta\), conclusion
\(\Gamma,\varphi\to\psi\Rightarrow\Delta\)): \(v\entails\varphi\to\psi\), so
it is not the case that \(v\entails\varphi\) and \(v\not\entails\psi\).  If
\(v\not\entails\varphi\), \(v\) falsifies the left premise; otherwise
\(v\entails\psi\) and \(v\) falsifies the right premise.

\(\to R\) (premise \(\Gamma,\varphi\Rightarrow\Delta,\psi\), conclusion
\(\Gamma\Rightarrow\Delta,\varphi\to\psi\)):
\(v\not\entails\varphi\to\psi\), so \(v\entails\varphi\) and
\(v\not\entails\psi\); \(v\) falsifies the premise.

Cut (premises \(\Gamma\Rightarrow\Delta,\varphi\) and
\(\Sigma,\varphi\Rightarrow\Pi\), conclusion
\(\Gamma,\Sigma\Rightarrow\Delta,\Pi\)): the members of
\(\operatorname{supp}(\Gamma)\cup\operatorname{supp}(\Sigma)\) are true and
those of \(\operatorname{supp}(\Delta)\cup\operatorname{supp}(\Pi)\) false.
If \(v\not\entails\varphi\), \(v\) falsifies the left premise; if
\(v\entails\varphi\), it falsifies the right premise.

\emph{Consequence.}  Show by induction on the height of a derivation that its
endsequent is valid.  An initial sequent is valid by the first paragraph.  If
the last rule had a conclusion falsified by some \(v\), then by the rule case
\(v\) would falsify a premise, which is valid by the induction hypothesis.
\end{proof}
\end{annexitem}

\begin{annexitem}{annex:p02-w02b-p02-009-nd-rule-soundness}{thm:p02-w02b-p02-009-nd-rule-soundness}{Every rule of natural deduction preserves truth from open assumptions}
\begin{proof}
For a derivation \(D\) let \(O(D)\) be the set of formulas having an open
occurrence in \(D\).  Prove by induction on \(D\): every valuation \(v\) with
\(v\entails\theta\) for all \(\theta\in O(D)\) satisfies the conclusion of
\(D\).  Write \(D_1,D_2,D_3\) for the immediate subderivations.

If the root rule discharges nothing, the open occurrences of \(D\) are
exactly those of its immediate subderivations, so \(O(D_m)\subseteq O(D)\)
and the induction hypothesis applies to every \(D_m\).  If the root rule
discharges occurrences of a formula \(\alpha\) above a premise \(D_m\), every
open occurrence of \(D_m\) is either open in \(D\) or a discharged occurrence
of \(\alpha\); hence a valuation satisfying \(O(D)\) and \(\alpha\) satisfies
\(O(D_m)\).  Fix \(v\) satisfying \(O(D)\).

\emph{Assumption} \(\varphi\): the occurrence is open, so
\(\varphi\in O(D)\) and \(v\entails\varphi\).

\(\land I\): \(v\entails\varphi\) and \(v\entails\psi\) by the induction
hypothesis, so \(v\entails\varphi\land\psi\).

\(\land E_1\), \(\land E_2\): \(v\entails\varphi\land\psi\), so
\(v\entails\varphi\) and \(v\entails\psi\).

\(\lor I_1\), \(\lor I_2\): from \(v\entails\varphi\), or from
\(v\entails\psi\), follows \(v\entails\varphi\lor\psi\).

\(\to E\): \(v\entails\varphi\) and \(v\entails\varphi\to\psi\).  If
\(v\not\entails\psi\), then \(v\entails\varphi\) and \(v\not\entails\psi\)
would give \(v\not\entails\varphi\to\psi\).  So \(v\entails\psi\).

\(\neg E\): the induction hypothesis would give \(v\entails\varphi\) and
\(v\entails\neg\varphi\), that is \(v\not\entails\varphi\).  So no
valuation satisfies \(O(D)\), and the claim holds for the conclusion
\(\bottom\).

\(\bottom E\): the induction hypothesis would give \(v\entails\bottom\),
which is impossible; again no valuation satisfies \(O(D)\).

\(\lor E\) with labels \(i,j\) and conclusion \(\chi\): no occurrence above
the first premise is discharged, so \(O(D_1)\subseteq O(D)\) and
\(v\entails\varphi\lor\psi\).  If \(v\entails\varphi\), then \(v\) satisfies
\(O(D_2)\), since the discharged occurrences above the second premise are
occurrences of \(\varphi\); hence \(v\entails\chi\).  Otherwise
\(v\entails\psi\), \(v\) satisfies \(O(D_3)\) and \(v\entails\chi\).

\(\to I\): if \(v\not\entails\varphi\), then \(v\entails\varphi\to\psi\).
If \(v\entails\varphi\), then \(v\) satisfies \(O(D_1)\), so
\(v\entails\psi\) and \(v\entails\varphi\to\psi\).

\(\neg I\): if \(v\entails\varphi\), then \(v\) satisfies \(O(D_1)\) and the
induction hypothesis gives \(v\entails\bottom\), impossible.  So
\(v\not\entails\varphi\) and \(v\entails\neg\varphi\).

\(\mathsf{RAA}\): the discharged occurrences are of \(\neg\varphi\).  If
\(v\not\entails\varphi\), then \(v\entails\neg\varphi\), so \(v\) satisfies
\(O(D_1)\) and \(v\entails\bottom\), impossible.  So \(v\entails\varphi\).

\emph{Consequence.}  If \(\Gamma\derives_{\mathsf{ND}_{\mathsf{CPL}}}\varphi\)
by a derivation \(D\), then \(O(D)\subseteq\Gamma\), so every \(v\) satisfying
\(\Gamma\) satisfies \(O(D)\) and hence \(\varphi\).
\end{proof}
\end{annexitem}

\begin{annexitem}{annex:p02-w02b-p02-009-inversion}{thm:p02-w02b-p02-009-inversion}{Semantic invertibility of the logical rules}
\begin{proof}
Let \(v\) falsify the indicated premise; then every member of its antecedent
support is true and every member of its succedent support false.

\(\neg L\): \(v\not\entails\varphi\) gives \(v\entails\neg\varphi\), so \(v\)
falsifies \(\Gamma,\neg\varphi\Rightarrow\Delta\).
\(\neg R\): \(v\entails\varphi\) gives \(v\not\entails\neg\varphi\), so \(v\)
falsifies \(\Gamma\Rightarrow\Delta,\neg\varphi\).
\(\land L\): \(v\entails\varphi\) and \(v\entails\psi\) give
\(v\entails\varphi\land\psi\).
\(\land R\): if \(v\) falsifies \(\Gamma\Rightarrow\Delta,\varphi\), then
\(v\not\entails\varphi\) gives \(v\not\entails\varphi\land\psi\); the right
premise is symmetric.
\(\lor L\): if \(v\) falsifies \(\Gamma,\varphi\Rightarrow\Delta\), then
\(v\entails\varphi\lor\psi\); the right premise is symmetric.
\(\lor R\): \(v\not\entails\varphi\) and \(v\not\entails\psi\) give
\(v\not\entails\varphi\lor\psi\).
\(\to L\): if \(v\) falsifies \(\Gamma\Rightarrow\Delta,\varphi\), then
\(v\not\entails\varphi\), so \(v\entails\varphi\to\psi\); if \(v\) falsifies
\(\Gamma,\psi\Rightarrow\Delta\), then \(v\entails\psi\), so again
\(v\entails\varphi\to\psi\).
\(\to R\): \(v\entails\varphi\) and \(v\not\entails\psi\) give
\(v\not\entails\varphi\to\psi\).
In each case all context formulas keep their values, so \(v\) falsifies the
conclusion.  If the conclusion is valid, no valuation falsifies it, so none
falsifies any premise.

\emph{Weight.}  Weight is additive:
\(\operatorname{wt}(\Gamma,\Pi\Rightarrow\Delta,\Sigma)=
\operatorname{wt}(\Gamma\Rightarrow\Delta)+\operatorname{wt}(\Pi\Rightarrow\Sigma)\),
and a single formula \(\chi\) on either side contributes \(c(\chi)\).  With
\(w\defeq\operatorname{wt}(\Gamma\Rightarrow\Delta)\) for the shared context,
the conclusion has weight \(w+c(\neg\varphi)=w+c(\varphi)+1\) in the
negation rules and \(w+c(\varphi)+c(\psi)+1\) in the binary rules.  The
premises have weight \(w+c(\varphi)\) (\(\neg L\), \(\neg R\), the left
premises of \(\land R\), \(\lor L\), \(\to L\)), \(w+c(\psi)\) (their right
premises), or \(w+c(\varphi)+c(\psi)\) (\(\land L\), \(\lor R\), \(\to R\)),
each strictly smaller.
\end{proof}
\end{annexitem}

\begin{annexitem}{annex:p02-w02b-p02-009-cutfree-completeness}{thm:p02-w02b-p02-009-cutfree-completeness}{Purely logical completeness}
\begin{proof}
By strong induction on \(n=\operatorname{wt}(\Gamma\Rightarrow\Delta)\), show
that every valid sequent of weight \(n\) has a purely logical derivation.
Let \(\Gamma\Rightarrow\Delta\) be valid of weight \(n\).

\emph{Case 1:} \(\bottom\in\operatorname{supp}(\Gamma)\).  Write
\(\Gamma=\Gamma',\bottom\), where \(\Gamma'\) is \(\Gamma\) with the
multiplicity of \(\bottom\) lowered by one.  The sequent is the
\(\bottom\)-sequent \(\Gamma',\bottom\Rightarrow\Delta\), a one-node
derivation.

\emph{Case 2:} some \(\varphi\) lies in both supports.  Write
\(\Gamma=\Gamma',\varphi\) and \(\Delta=\varphi,\Delta'\) in the same way.
The sequent is the identity sequent
\(\Gamma',\varphi\Rightarrow\varphi,\Delta'\).

\emph{Case 3:} neither case applies and every member of
\(\operatorname{supp}(\Gamma)\cup\operatorname{supp}(\Delta)\) is a variable
or \(\bottom\).  Let \(v(p_n)=1\) iff \(p_n\in\operatorname{supp}(\Gamma)\).
Every member of \(\operatorname{supp}(\Gamma)\) is a variable, since
\(\bottom\notin\operatorname{supp}(\Gamma)\), and is true under \(v\).  Every
member of \(\operatorname{supp}(\Delta)\) is \(\bottom\), false under \(v\),
or a variable not in \(\operatorname{supp}(\Gamma)\), since the supports are
disjoint, and hence false under \(v\).  So \(v\) falsifies the sequent,
contrary to validity; this case does not occur.

\emph{Case 4:} neither Case 1 nor Case 2 applies and some
\(\chi\in\operatorname{supp}(\Gamma)\cup\operatorname{supp}(\Delta)\) is of
the form \(\neg\varphi\), \(\varphi\land\psi\), \(\varphi\lor\psi\) or
\(\varphi\to\psi\); by unique readability
(\cref{thm:prop-unique-readability}) exactly one form applies, with unique
\(\varphi,\psi\).  If \(\chi\in\operatorname{supp}(\Gamma)\), write
\(\Gamma=\Gamma',\chi\); then \(\Gamma\Rightarrow\Delta\) is the conclusion of
the instance of \(\neg L\), \(\land L\), \(\lor L\) or \(\to L\) with context
\(\Gamma',\Delta\) and principal formula \(\chi\).  If
\(\chi\in\operatorname{supp}(\Delta)\), write \(\Delta=\Delta',\chi\) and use
\(\neg R\), \(\land R\), \(\lor R\) or \(\to R\) with context
\(\Gamma,\Delta'\).  By \cref{thm:p02-w02b-p02-009-inversion}, each premise of
that instance is valid and has weight less than \(n\); by the induction
hypothesis each has a purely logical derivation.  Placing the rule beneath
them gives a purely logical derivation of \(\Gamma\Rightarrow\Delta\).

The four cases are exhaustive: if Cases 1, 2 and 4 fail, every formula in the
supports is a variable or \(\bottom\), which is Case 3.  No step uses Cut,
weakening or contraction.
\end{proof}
\end{annexitem}

\begin{annexitem}{annex:p02-w02b-p02-009-cut-admissibility}{thm:p02-w02b-p02-009-cut-admissibility}{Cut admissibility for the classical propositional sequent calculus}
\begin{proof}
If \(\Gamma\Rightarrow\Delta\) is derivable, it is valid by
\cref{thm:p02-w02b-p02-009-lk-rule-soundness}, whose rule cases include Cut,
weakening and contraction.  By
\cref{thm:p02-w02b-p02-009-cutfree-completeness} it has a purely logical
derivation, which contains no Cut, weakening or contraction.  This proves
the first assertion and both admissibility statements.  For the
equivalences: a purely logical derivation is a derivation; a derivable
sequent is valid by the soundness lemma; a valid sequent has a purely logical
derivation by the completeness theorem.
\end{proof}
\end{annexitem}

\begin{annexitem}{annex:p02-w02b-p02-009-subformula}{thm:p02-w02b-p02-009-subformula}{Subformula property without Cut}
\begin{proof}
\emph{Facts about \(\operatorname{Sub}\).}  (a) \(\chi\in\operatorname{Sub}(\chi)\)
for every \(\chi\), since every clause of the definition in
\cref{sec:prop-syntax-semantics-01} puts \(\chi\) into
\(\operatorname{Sub}(\chi)\).  (b) \(\varphi\in\operatorname{Sub}(\neg\varphi)\)
and \(\varphi,\psi\in\operatorname{Sub}(\varphi\circ\psi)\), by (a) and the
clauses for \(\neg\) and \(\circ\).  (c) If \(\psi\in\operatorname{Sub}(\chi)\),
then \(\operatorname{Sub}(\psi)\subseteq\operatorname{Sub}(\chi)\).  Prove (c)
by induction on \(\chi\).  If \(\psi=\chi\) there is nothing to prove; this
covers \(\chi=p_n\) and \(\chi=\bottom\), whose subformula set is
\(\{\chi\}\).  If \(\chi=\neg\varphi\) and \(\psi\neq\chi\), then
\(\psi\in\operatorname{Sub}(\varphi)\), so by the induction hypothesis
\(\operatorname{Sub}(\psi)\subseteq\operatorname{Sub}(\varphi)
\subseteq\operatorname{Sub}(\chi)\).  If \(\chi=\varphi\circ\psi'\) and
\(\psi\neq\chi\), then \(\psi\) lies in \(\operatorname{Sub}(\varphi)\) or in
\(\operatorname{Sub}(\psi')\), and the same argument applies.

\emph{Rule step.}  Consider an application of a rule other than Cut, with
conclusion \(S\), and let \(\theta\) lie in the support of a side of a
premise.  For \(W_L\) and \(W_R\) each premise support is included in the
corresponding conclusion support; for \(C_L\) and \(C_R\) the supports are
equal; so \(\theta\) is itself a formula of \(S\).  For a logical rule, the
support of each premise side is the support of the corresponding context of
\(S\), possibly together with one or two side formulas.  A context formula is
a formula of \(S\).  A side formula is \(\varphi\) or \(\psi\), and the
principal formula of \(S\) is \(\neg\varphi\), \(\varphi\land\psi\),
\(\varphi\lor\psi\) or \(\varphi\to\psi\); by (b) the side formula is a
subformula of a formula of \(S\).  With (a), in every case \(\theta\in
\operatorname{Sub}(\chi')\) for some formula \(\chi'\) of \(S\).

\emph{Induction.}  Let \(D\) be Cut-free with endsequent
\(\Gamma\Rightarrow\Delta\).  By induction on the number of rule applications
between a node and the root, every formula \(\theta\) of the sequent at that
node lies in \(\operatorname{Sub}(\chi)\) for some
\(\chi\in\operatorname{supp}(\Gamma)\cup\operatorname{supp}(\Delta)\).  At
the root this is (a).  At a premise of a rule application whose conclusion
node satisfies the claim, the rule step gives \(\theta\in
\operatorname{Sub}(\chi')\) for a formula \(\chi'\) of the conclusion, and
\(\chi'\in\operatorname{Sub}(\chi)\) for an endsequent formula \(\chi\); by
(c), \(\theta\in\operatorname{Sub}(\chi)\).
\end{proof}
\end{annexitem}

\begin{annexitem}{annex:p02-w02b-p02-009-analytic}{thm:p02-w02b-p02-009-analytic}{Analytic derivations}
\begin{proof}
A derivable sequent has a purely logical derivation by
\cref{thm:p02-w02b-p02-009-cut-admissibility}; it contains no Cut, so it has
the subformula property by \cref{thm:p02-w02b-p02-009-subformula}.  Each
\(\operatorname{Sub}(\chi)\) is finite, by induction on \(\chi\) from its
recursive clauses, and the supports are finite, so the displayed union is
finite.  Every valuation falsifies the empty sequent, since both supports are
empty; it is not valid, and by
\cref{thm:p02-w02b-p02-009-lk-rule-soundness} it is not derivable.
\end{proof}
\end{annexitem}

\begin{annexitem}{annex:p02-w02b-p02-009-nd-to-lk}{thm:p02-w02b-p02-009-nd-to-lk}{Natural deduction into the sequent calculus}
\begin{proof}
By induction on \(D\), check that each tree of
\cref{def:p02-w02b-p02-009-nd-to-lk} consists of instances of the rules of
\(\mathsf{LK}_{\mathsf{CPL}}\) and has endsequent
\(\operatorname{om}(D)\Rightarrow\varphi\).  By the induction hypothesis the
premises \(\Gamma_m\Rightarrow\cdot\) are derivable with
\(\Gamma_m=\operatorname{om}(D_m)\).  A rule instance is named by its
schematic contexts.

\emph{Assumption.}  \(D\) is one open occurrence of \(\varphi\), so
\(\operatorname{om}(D)=\varphi\); \(\varphi\Rightarrow\varphi\) is the
identity sequent with empty contexts.

\emph{Rules that discharge nothing.}  Then
\(\operatorname{om}(D)=\Gamma_1+\dots+\Gamma_n\) by
\cref{def:p02-w02b-p02-009-open-multiset}.
\(\land I\): \(|\Gamma_2|\) applications of \(W_L\) turn
\(\Gamma_1\Rightarrow\varphi\) into \(\Gamma_1,\Gamma_2\Rightarrow\varphi\),
and \(|\Gamma_1|\) applications turn \(\Gamma_2\Rightarrow\psi\) into
\(\Gamma_1,\Gamma_2\Rightarrow\psi\); \(\land R\) with context
\(\Gamma_1,\Gamma_2\) and empty \(\Delta\) gives
\(\Gamma_1,\Gamma_2\Rightarrow\varphi\land\psi\).
\(\land E_1\): \(\varphi,\psi\Rightarrow\varphi\) is an identity sequent
(antecedent context \(\psi\)); \(\land L\) with empty contexts gives
\(\varphi\land\psi\Rightarrow\varphi\); Cut with \(\Gamma=\Gamma_1\),
\(\Delta,\Sigma\) empty and \(\Pi=\varphi\) gives
\(\Gamma_1\Rightarrow\varphi\).  \(\land E_2\) is the same with
\(\varphi,\psi\Rightarrow\psi\).
\(\lor I_1\): \(W_R\) gives \(\Gamma_1\Rightarrow\varphi,\psi\) and
\(\lor R\) with empty \(\Delta\) gives
\(\Gamma_1\Rightarrow\varphi\lor\psi\); \(\lor I_2\) is symmetric.
\(\to E\): \(W_R\) gives \(\Gamma_1\Rightarrow\psi,\varphi\);
\(\Gamma_1,\psi\Rightarrow\psi\) is an identity sequent; \(\to L\) with
context \(\Gamma_1\) and \(\Delta=\psi\) gives
\(\Gamma_1,\varphi\to\psi\Rightarrow\psi\); Cut with left premise
\(\Gamma_2\Rightarrow\varphi\to\psi\) (\(\Delta\) empty) and \(\Sigma=\Gamma_1\),
\(\Pi=\psi\) gives \(\Gamma_2,\Gamma_1\Rightarrow\psi\).
\(\neg E\): \(\neg L\) with context \(\Gamma_1\) and empty \(\Delta\) gives
\(\Gamma_1,\neg\varphi\Rightarrow{}\); \(W_R\) gives
\(\Gamma_1,\neg\varphi\Rightarrow\bottom\); Cut with left premise
\(\Gamma_2\Rightarrow\neg\varphi\) gives \(\Gamma_2,\Gamma_1\Rightarrow\bottom\).
\(\bottom E\): \(\bottom\Rightarrow\varphi\) is a \(\bottom\)-sequent; Cut
with left premise \(\Gamma_1\Rightarrow\bottom\) gives
\(\Gamma_1\Rightarrow\varphi\).
In each case the antecedent is \(\Gamma_1+\dots+\Gamma_n\), since multiset
sum is commutative.

\emph{\(\to I\), \(\neg I\), \(\mathsf{RAA}\) with label \(i\).}  The root
discharges exactly the \(k\) open occurrences in \(D_1\) of the bracketed
formula \(\alpha\) (\(\alpha=\varphi\), \(\varphi\), \(\neg\varphi\)) marked
\(i\), and \(k\geq1\) by the nonvacuity clause.  All other open occurrences of
\(D_1\), including occurrences of \(\alpha\) with another mark, stay open.
So \(\Gamma_1=\operatorname{om}(D)+\alpha^{k}\); write
\(\Gamma'=\operatorname{om}(D)\).  Since \(k\geq1\), \(k-1\) applications of
\(C_L\) to \(\alpha\) turn \(\Gamma',\alpha^{k}\Rightarrow\cdot\) into
\(\Gamma',\alpha\Rightarrow\cdot\).  For \(\to I\), \(\to R\) with context
\(\Gamma'\) gives \(\Gamma'\Rightarrow\varphi\to\psi\).  For \(\neg I\), Cut
with left premise \(\Gamma',\varphi\Rightarrow\bottom\) and right premise the
\(\bottom\)-sequent \(\bottom\Rightarrow{}\) (\(\Sigma,\Pi\) empty) gives
\(\Gamma',\varphi\Rightarrow{}\), and \(\neg R\) gives
\(\Gamma'\Rightarrow\neg\varphi\).  For \(\mathsf{RAA}\), the same Cut gives
\(\Gamma',\neg\varphi\Rightarrow{}\); \(\neg R\) applied to the identity
sequent \(\varphi\Rightarrow\varphi\), with empty antecedent context and
\(\Delta=\varphi\), gives \({}\Rightarrow\varphi,\neg\varphi\); Cut on
\(\neg\varphi\) with \(\Gamma\) empty, \(\Delta=\varphi\), \(\Sigma=\Gamma'\)
and \(\Pi\) empty gives \(\Gamma'\Rightarrow\varphi\).

\emph{\(\lor E\) with labels \(i,j\).}  The root discharges the \(k\geq1\)
open occurrences of \(\varphi\) marked \(i\) in \(D_2\) and the \(m\geq1\)
open occurrences of \(\psi\) marked \(j\) in \(D_3\), and nothing in \(D_1\).
Hence \(\operatorname{om}(D)=\Gamma_1+\Gamma_2'+\Gamma_3'\) with
\(\Gamma_2=\Gamma_2'+\varphi^{k}\) and \(\Gamma_3=\Gamma_3'+\psi^{m}\).
\(k-1\) applications of \(C_L\) and \(|\Gamma_3'|\) of \(W_L\) give
\(\Gamma_2',\Gamma_3',\varphi\Rightarrow\chi\); \(m-1\) of \(C_L\) and
\(|\Gamma_2'|\) of \(W_L\) give \(\Gamma_2',\Gamma_3',\psi\Rightarrow\chi\).
\(\lor L\) with context \(\Gamma_2',\Gamma_3'\) and \(\Delta=\chi\) gives
\(\Gamma_2',\Gamma_3',\varphi\lor\psi\Rightarrow\chi\), and Cut with left
premise \(\Gamma_1\Rightarrow\varphi\lor\psi\) gives
\(\Gamma_1,\Gamma_2',\Gamma_3'\Rightarrow\chi\).

The counts of \(C_L\) and \(W_L\) applications stated in the theorem are
read off from these cases; no other case uses \(C_L\) or \(W_L\).
\end{proof}
\end{annexitem}

\begin{annexitem}{annex:p02-w02b-p02-009-lk-to-nd}{thm:p02-w02b-p02-009-lk-to-nd}{Sequent derivations into natural deduction}
\begin{proof}
Call a derivation with conclusion \(\bottom\) \emph{well marked} iff each
of its open occurrences is marked by a value of \(\ell_{\mathrm L}\),
\(\ell_{\mathrm R}\) or \(\ell_{\mathrm A}\), and an occurrence marked
\(\ell_{\mathrm L}(\chi)\) is an occurrence of \(\chi\), one marked
\(\ell_{\mathrm R}(\chi)\) is an occurrence of \(\neg\chi\), and one marked
\(\ell_{\mathrm A}(\chi)\) is an occurrence of \(\chi\).  Joint injectivity
makes the mark of an occurrence determine its kind and its formula
\(\chi\).  A derivation representing a sequent is well marked and has no
occurrence marked by a value of \(\ell_{\mathrm A}\).

\emph{Helper facts.}  Let \(D\) be well marked.
(H1) \(\operatorname{ext}_\chi(D)\) is a derivation of \(\chi\) whose open
occurrences are those of \(D\) other than the ones marked
\(\ell_{\mathrm R}(\chi)\).  Indeed, \(\mathsf{RAA}\) with label
\(\ell_{\mathrm R}(\chi)\) discharges the open occurrences of \(\neg\chi\)
marked \(\ell_{\mathrm R}(\chi)\), and in a well-marked derivation every open
occurrence with that mark is an occurrence of \(\neg\chi\).  The application
is made only when such an occurrence exists, so it is nonvacuous; otherwise
\(\bottom E\) discharges nothing and no occurrence marked
\(\ell_{\mathrm R}(\chi)\) was open.
(H2) In the same way \(\operatorname{neg}_\chi(D)\) is a derivation of
\(\neg\chi\) whose open occurrences are those of \(D\) other than the ones
marked \(\ell_{\mathrm L}(\chi)\).
(H3) \(\operatorname{pin}_\chi(D)\) is a well-marked derivation of
\(\bottom\) whose open occurrences are those of
\(\operatorname{neg}_\chi(D)\) together with the new leaf, an occurrence of
\(\chi\) marked \(\ell_{\mathrm A}(\chi)\); the \(\neg E\) step has premises
\(\chi\) and \(\neg\chi\) and discharges nothing.  If \(D\) had no open
occurrence marked by a value of \(\ell_{\mathrm A}\), the new leaf is the only
such occurrence of \(\operatorname{pin}_\chi(D)\).
(H4) A tree obtained by applying \(\neg E\), \(\land I\), \(\land E_m\),
\(\lor I_m\) or \(\to E\) to derivations and to new marked leaves has as open
occurrences exactly the open occurrences of its parts.

Now induct on the \(\mathsf{LK}_{\mathsf{CPL}}\)-derivation.  In each case the
displayed tree is built from rule instances of
\(\mathsf{ND}_{\mathsf{CPL}}\) whose premises and conclusions have the
displayed forms; it remains to list the open occurrences and check that each
is a left occurrence of a member of the conclusion's antecedent support or a
right occurrence of a member of its succedent support.  By the induction
hypothesis \(D_1\) (and \(D_2\)) represent the premises.

Identity \(\Gamma,\varphi\Rightarrow\varphi,\Delta\): the open occurrences
are \(\varphi^{\ell_{\mathrm L}(\varphi)}\), a left occurrence with
\(\varphi\in\operatorname{supp}(\Gamma,\varphi)\), and
\((\neg\varphi)^{\ell_{\mathrm R}(\varphi)}\), a right occurrence with
\(\varphi\in\operatorname{supp}(\varphi,\Delta)\).  \(\bottom\)-sequent: the
single occurrence \(\bottom^{\ell_{\mathrm L}(\bottom)}\) is a left occurrence
and \(\bottom\in\operatorname{supp}(\Gamma,\bottom)\).

\(W_L,W_R\): the supports of the conclusion include those of the premise, so
\(D_1\) represents the conclusion.  \(C_L,C_R\): the supports are equal.

\(\neg L\): by (H1) the open occurrences of
\(\operatorname{ext}_\varphi(D_1)\) are left occurrences of members of
\(\operatorname{supp}(\Gamma)\) and right occurrences of members of
\(\operatorname{supp}(\Delta,\varphi)\setminus\{\varphi\}\subseteq
\operatorname{supp}(\Delta)\).  The new leaf
\((\neg\varphi)^{\ell_{\mathrm L}(\neg\varphi)}\) is a left occurrence of
\(\neg\varphi\in\operatorname{supp}(\Gamma,\neg\varphi)\).  By (H4) these are
all open occurrences.

\(\neg R\): by (H2), \(\operatorname{neg}_\varphi(D_1)\) has left occurrences
of members of \(\operatorname{supp}(\Gamma,\varphi)\setminus\{\varphi\}
\subseteq\operatorname{supp}(\Gamma)\) and right occurrences of members of
\(\operatorname{supp}(\Delta)\).  The leaf
\((\neg\neg\varphi)^{\ell_{\mathrm R}(\neg\varphi)}\) is a right occurrence of
\(\neg\varphi\in\operatorname{supp}(\Delta,\neg\varphi)\).

\(\land L\): by (H2) and (H4), \(D'\) has left occurrences of members of
\(\operatorname{supp}(\Gamma,\varphi,\psi)\setminus\{\psi\}\), the left
occurrence \((\varphi\land\psi)^{\ell_{\mathrm L}(\varphi\land\psi)}\), and
right occurrences of members of \(\operatorname{supp}(\Delta)\); it is well
marked and represents \(\Gamma,\varphi,\varphi\land\psi\Rightarrow\Delta\).
Applying (H2) to \(D'\) removes exactly the occurrences marked
\(\ell_{\mathrm L}(\varphi)\); the leaf marked
\(\ell_{\mathrm L}(\varphi\land\psi)\) is not among them, because
\(\varphi\land\psi\neq\varphi\).  The final tree therefore has left
occurrences of members of \(\operatorname{supp}(\Gamma)\) and of
\(\varphi\land\psi\), and right occurrences of members of
\(\operatorname{supp}(\Delta)\).  If \(\varphi=\psi\), the first
\(\operatorname{neg}\) has already removed the occurrences of \(\varphi\) and
the second uses \(\bottom E\); the conclusion is the same.

\(\land R\): by (H1) the open occurrences of
\(\operatorname{ext}_\varphi(D_1)\) and \(\operatorname{ext}_\psi(D_2)\) are
left occurrences of members of \(\operatorname{supp}(\Gamma)\) and right
occurrences of members of \(\operatorname{supp}(\Delta)\); the leaf
\((\neg(\varphi\land\psi))^{\ell_{\mathrm R}(\varphi\land\psi)}\) is a right
occurrence of \(\varphi\land\psi\).

\(\lor L\): by (H3), \(\operatorname{pin}_\varphi(D_1)\) has left occurrences
of members of \(\operatorname{supp}(\Gamma)\), right occurrences of members of
\(\operatorname{supp}(\Delta)\), and exactly one occurrence marked by a value
of \(\ell_{\mathrm A}\), namely \(\varphi^{\ell_{\mathrm A}(\varphi)}\);
likewise for \(\operatorname{pin}_\psi(D_2)\) with
\(\psi^{\ell_{\mathrm A}(\psi)}\).  The \(\lor E\) step with labels
\(\ell_{\mathrm A}(\varphi)\) and \(\ell_{\mathrm A}(\psi)\) discharges exactly
these two occurrences, one above each minor premise, and nothing above the
major premise, the leaf \((\varphi\lor\psi)^{\ell_{\mathrm L}(\varphi\lor\psi)}\);
each label discharges one occurrence.  The remaining open occurrences are
left occurrences of members of \(\operatorname{supp}(\Gamma)\) and of
\(\varphi\lor\psi\), and right occurrences of members of
\(\operatorname{supp}(\Delta)\).

\(\lor R\): by (H1), \(\operatorname{ext}_\psi(D_1)\) has left occurrences of
members of \(\operatorname{supp}(\Gamma)\) and right occurrences of members of
\(\operatorname{supp}(\Delta,\varphi,\psi)\setminus\{\psi\}\); so \(D'\)
represents \(\Gamma\Rightarrow\Delta,\varphi,\varphi\lor\psi\).  By (H1)
again, \(\operatorname{ext}_\varphi(D')\) removes exactly the occurrences
marked \(\ell_{\mathrm R}(\varphi)\); the leaf marked
\(\ell_{\mathrm R}(\varphi\lor\psi)\) stays, because
\(\varphi\lor\psi\neq\varphi\).  The final tree has left occurrences of
members of \(\operatorname{supp}(\Gamma)\) and right occurrences of members
of \(\operatorname{supp}(\Delta)\cup\{\varphi\lor\psi\}\).

\(\to L\): by (H1) and (H2), \(\operatorname{ext}_\varphi(D_1)\) and
\(\operatorname{neg}_\psi(D_2)\) have left occurrences of members of
\(\operatorname{supp}(\Gamma)\) and right occurrences of members of
\(\operatorname{supp}(\Delta)\).  The \(\to E\) step has premises
\(\varphi\) and \(\varphi\to\psi\), the latter the left occurrence
\((\varphi\to\psi)^{\ell_{\mathrm L}(\varphi\to\psi)}\); the \(\neg E\) step
has premises \(\psi\) and \(\neg\psi\).

\(\to R\): by (H3), \(\operatorname{pin}_\varphi(D_1)\) has left occurrences
of members of \(\operatorname{supp}(\Gamma)\), right occurrences of members
of \(\operatorname{supp}(\Delta,\psi)\), and the single occurrence
\(\varphi^{\ell_{\mathrm A}(\varphi)}\).  By (H1),
\(\operatorname{ext}_\psi\) removes the right occurrences of \(\psi\) and
keeps the occurrence marked \(\ell_{\mathrm A}(\varphi)\), whose mark differs
from \(\ell_{\mathrm R}(\psi)\).  The \(\to I\) step with label
\(\ell_{\mathrm A}(\varphi)\) discharges exactly that occurrence of
\(\varphi\), so it is nonvacuous and leaves left occurrences of members of
\(\operatorname{supp}(\Gamma)\) and right occurrences of members of
\(\operatorname{supp}(\Delta)\).  The leaf
\((\neg(\varphi\to\psi))^{\ell_{\mathrm R}(\varphi\to\psi)}\) is a right
occurrence of \(\varphi\to\psi\in\operatorname{supp}(\Delta,\varphi\to\psi)\).

Cut: by (H1), \(\operatorname{ext}_\varphi(D_1)\) has left occurrences of
members of \(\operatorname{supp}(\Gamma)\) and right occurrences of members
of \(\operatorname{supp}(\Delta)\); by (H2), \(\operatorname{neg}_\varphi(D_2)\)
has left occurrences of members of \(\operatorname{supp}(\Sigma)\) and right
occurrences of members of \(\operatorname{supp}(\Pi)\).  The \(\neg E\) step
has premises \(\varphi\) and \(\neg\varphi\).  All these supports lie in those
of \(\Gamma,\Sigma\Rightarrow\Delta,\Pi\).

In every case the result has conclusion \(\bottom\) and only left and right
occurrences as open occurrences, so it represents the conclusion sequent.
The statements about discharged occurrences are (H1)--(H3) and the
\(\lor L\) and \(\to R\) cases; the statement about weakening and
contraction is the \(W\) and \(C\) cases.
\end{proof}
\end{annexitem}

\begin{annexitem}{annex:p02-w02b-p02-009-equivalence}{thm:p02-w02b-p02-009-equivalence}{Derivability correspondence by explicit translation}
\begin{proof}
(1) Let \(D\) derive \(\varphi\) with every open assumption in \(\Gamma\).  By
\cref{thm:p02-w02b-p02-009-nd-to-lk},
\(\operatorname{om}(D)\Rightarrow\varphi\) is derivable, and
\(\operatorname{supp}(\operatorname{om}(D))\subseteq\Gamma\).  For each
\(\chi\) in this support apply \(C_L\) \(\operatorname{om}(D)(\chi)-1\)
times; the antecedent becomes the multiset with multiplicity one on
\(\operatorname{supp}(\operatorname{om}(D))\).  Then apply \(W_L\) once for
each member of \(\Gamma\setminus\operatorname{supp}(\operatorname{om}(D))\);
the antecedent becomes \(\Gamma\) with multiplicity one.  A purely logical
derivation exists by \cref{thm:p02-w02b-p02-009-cut-admissibility}.

(2) By \cref{thm:p02-w02b-p02-009-lk-to-nd} some \(D\) represents
\(\Gamma\Rightarrow\Delta\): it concludes \(\bottom\), and each open
assumption is a member of \(\Gamma\) or \(\neg\delta\) with
\(\delta\in\operatorname{supp}(\Delta)\).  This is the first claim.  If
\(\Delta=\varphi\), every right occurrence of \(D\) is marked
\(\ell_{\mathrm R}(\varphi)\); by (H1) in
\cref{annex:p02-w02b-p02-009-lk-to-nd}, \(\operatorname{ext}_\varphi(D)\)
derives \(\varphi\) and its open occurrences are the left occurrences of
\(D\), all of members of \(\Gamma\).

Only \cref{thm:p02-w02b-p02-009-nd-to-lk,thm:p02-w02b-p02-009-lk-to-nd} and
structural rules are used, apart from the final sentence of (1).
\end{proof}
\end{annexitem}
